
\documentclass[11pt,a4paper]{article}
\usepackage[utf8]{inputenc}
\usepackage[margin=1.2in]{geometry}
\usepackage{amsmath}
\usepackage{amssymb}
\usepackage{listings}
\usepackage{xcolor}
\usepackage{hyperref}
\usepackage{titlesec}

% Professional styling
\definecolor{accentblue}{RGB}{30, 80, 150}
\definecolor{codegray}{rgb}{0.95,0.95,0.95}
\definecolor{codeblue}{RGB}{0, 0, 128}

\titleformat{\section}{\large\bfseries\color{accentblue}}{}{0em}{}[\titlerule]
\titleformat{\subsection}{\normalsize\bfseries}{}{0em}{}

\lstset{
    backgroundcolor=\color{codegray},
    basicstyle=\ttfamily\small,
    breaklines=true,
    frame=single,
    rulecolor=\color{accentblue!30},
    commentstyle=\color{gray},
    keywordstyle=\color{codeblue},
    showstringspaces=false
}

\title{
    \textbf{Repository Summary: Eternal Torment}\\
    \large \texttt{pale-shadow/eternal-torment}
}
\author{Franklin Edward Diaz}
\date{May 5, 2026}

\begin{document}

\maketitle

\section{Overview}
\textbf{Purpose:} A high-entropy sandbox repository dedicated to low-level coding practice, infrastructure-as-code (IaC) modeling, and continuous learning.

\textbf{Context:} Part of the \texttt{pale-shadow} organization, this repository serves as a functional testing ground for CTF tools, assembly experiments, and the recently integrated \texttt{ml-gnn} machine learning project. It bridges the gap between raw hardware interaction and cloud-scale infrastructure analysis.

\section{Recent Updates: Project \texttt{ml-gnn}}
Recent activity focuses on the intersection of Graph Neural Networks and Cloud Security. The primary objective is the automated extraction and analysis of Terraform-based infrastructure digraphs.

\begin{itemize}
    \item \textbf{Infrastructure Graph Collection:} A Python pipeline (\texttt{src/collection/}) that extracts relationship data from Terraform plans.
    \item \textbf{Data Lifecycle Management:} Integration of DVC (Data Version Control) for dataset tracking and \texttt{.metadata.json} consistency.
    \item \textbf{GCP Integration:} Automated exfiltration of DOT, PNG, and Terraform state files to the \texttt{backend-datastore} storage bucket.
\end{itemize}

\section{Setup \& Environment Configuration}
The project now supports multi-environment builds via GNU Autotools and containerized runtimes.

\subsection{Docker Workflow}
For isolated collection without managing local dependencies:
\begin{lstlisting}[language=bash]
docker build -t gnn-collection .
\end{lstlisting}

\subsection{Native Environment (Bash)}
Standard setup for Debian and OpenBSD environments:
\begin{lstlisting}[language=bash]
./bootstrap.sh
./configure
make python
export BASE_DIR=$(pwd)
source ${BASE_DIR}/_build/bin/activate
export PYTHONPATH=${BASE_DIR}
\end{lstlisting}

\subsection{Native Environment (Fish Shell)}
Setup for interactive Fish shell sessions:
\begin{lstlisting}[language=bash]
./bootstrap.sh
./configure
make python
source /home/franklin/workspace/ml-gnn/_build/bin/activate.fish
set --export PYTHONPATH /home/franklin/workspace/ml-gnn
\end{lstlisting}

\section{Execution}
With the environment active, the collection engine is executed from within a Terraform workspace directory:

\begin{lstlisting}[language=bash]
cd /path/to/terraform/workspace
python3 ${BASE_DIR}/src/collection/collection.py
\end{lstlisting}

\section{Computational Performance Metrics}
Numerical processing for GNN embeddings and large-scale graph analysis is offloaded to the \texttt{lab.bitsmasher.net} NVIDIA Jetson cluster. Throughput is modeled as:

\begin{equation}
P_{\text{total}} = \sum_{i=1}^{n} (N_{\text{cores}, i} \times f_{i} \times \text{FLOPS}_{\text{cycle}})
\end{equation}

Where $n=4$ nodes, $N_{\text{cores}, i}$ represents CUDA cores, and $f_{i}$ represents clock frequency.

\section{Technical Stack}
\begin{itemize}
    \item \textbf{Languages:} Assembly (AArch64), C, Python, Shell (Bash/Fish/Ksh).
    \item \textbf{Build Systems:} GNU Autotools, Make.
    \item \textbf{Cloud/IaC:} Terraform, Google Cloud Platform (GKE, Storage).
    \item \textbf{Security:} Bandit (SAST), Tox (Automation).
\end{itemize}

\vspace{2em}
\noindent\textit{Maintainer: devsecfranklin}\\
\textit{Organization: Pale Shadow}\\
\textit{License: Apache-2.0}

\end{document}
